\subsection{\texorpdfstring{第~\ref{sec:nonholonomic-constraints}~节}{非完整约束}}\label{subsec:solutions-nonholonomic-constraints}

\begin{description}[font=\normalfont,leftmargin=2em,labelwidth=1.5em]
\RGsolution{ex:integrability-implies-bracket-closure}{integrability-implies-bracket-closure}
设 $\Sigma$ 是可积分布. 那么对于每个 $p\in M$, 都存在 $p$ 的开邻域 $U$ 和局部坐标 $(x^1,\ldots,x^n):U\to\mathbb{R}^n$, 使得各叶与 $U$ 的交集的连通分支是 $(x^{m+1},\ldots,x^n):U\to\mathbb{R}^{n-m}$ 的水平集. 因此若 $X,Y\in\mathfrak X(\Sigma)$, 则在 $U$ 上
\[
  X=\sum_{i=1}^m X^i\frac{\partial}{\partial x^i},\quad
  Y=\sum_{i=1}^m Y^i\frac{\partial}{\partial x^i}.
\]
从而在 $U$ 上
\[
  [X,Y]=\sum_{i=1}^m(X\cdot Y^i-Y\cdot X^i)\frac{\partial}{\partial x^i}.
\]
因为 $U$ 是任取的, 所以可知 $[X,Y]\in\mathfrak X(U)$.

\RGsolution{ex:prove-exterior-forms-integrability}{prove-exterior-forms-integrability}
在局部总可以把 $\{\omega^1,\ldots,\omega^{n-m}\}$ 补成局部余标架
\[
  \{\theta^1,\ldots,\theta^m,\omega^1,\ldots,\omega^{n-m}\}.
\]
设 $\{X_1,\ldots,X_m,Y_1,\ldots,Y_{n-m}\}$ 是对偶局部标架, 则 $\Sigma$ 在局部由 $\{X_1,\ldots,X_m\}$ 给出. 有
\begin{align*}
  \mathrm{d}\omega^i(X_j,X_k)
  &=X_j\cdot\omega^i(X_k)-X_k\cdot\omega^i(X_j)-\omega^i([X_j,X_k])\\
  &=-\omega^i([X_j,X_k]),
\end{align*}
所以分布可积, 当且仅当对所有 $i=1,\ldots,n-m$ 和 $j,k=1,\ldots,m$, 有
\[
  \mathrm{d}\omega^i(X_j,X_k)=0.
\]
写作
\begin{align*}
  \mathrm{d}\omega^i
  &=\sum_{j,k=1}^m a^i_{jk}\theta^j\wedge\theta^k
    +\sum_{j=1}^m\sum_{k=1}^{n-m}b^i_{jk}\theta^j\wedge\omega^k\\
  &\quad+\sum_{j,k=1}^{n-m}c^i_{jk}\omega^j\wedge\omega^k,
\end{align*}
则
\[
  \mathrm{d}\omega^i(X_j,X_k)=a^i_{jk},
\]
因而可积性条件等价于要求各函数 $a^i_{jk}$ 为零. 由于
\[
  \mathrm{d}\omega^i\wedge\omega^1\wedge\cdots\wedge\omega^{n-m}
    =\sum_{j,k=1}^m a^i_{jk}\theta^j\wedge\theta^k\wedge\omega^1\wedge\cdots\wedge\omega^{n-m},
\]
这也就是要求
\[
  \mathrm{d}\omega^i\wedge\omega^1\wedge\cdots\wedge\omega^{n-m}=0
\]
(假定 $m\ge2$; 当 $m=1$ 时, 这个条件和可积性条件都显然满足).

\RGsolution{ex:prove-nonholonomic-energy-conservation}{prove-nonholonomic-energy-conservation}
有
\begin{align*}
  \frac{\mathrm{d}E}{\mathrm{d}t}
  &=\frac{\mathrm{d}}{\mathrm{d}t}\left(\frac12\langle\dot c(t),\dot c(t)\rangle+U(c(t))\right)\\
  &=\left\langle\frac{D\dot c}{\mathrm{d}t}(t),\dot c(t)\right\rangle
      +(\mathrm{d}U)_{c(t)}(\dot c(t))\\
  &=\mu\left(\frac{D\dot c}{\mathrm{d}t}\right)(\dot c)-\mathcal F(\dot c)(\dot c)
    =\mathcal R(\dot c)(\dot c)=0,
\end{align*}
因为对理想约束反力, 有
\[
  \mathcal R(\dot c)(\dot c)=\left\langle\mu^{-1}(\mathcal R(\dot c)),\dot c\right\rangle=0.
\]

\RGsolution{ex:ice-skate-dynamics}{ice-skate-dynamics}
\begin{enumerate}[label=(\alph*),ref=\alph*,leftmargin=*]
\item 给定 $\mathbb{R}^2\times S^1$ 中两点 $p=(x_0,y_0,\theta_0)$ 和 $q=(x_1,y_1,\theta_1)$, 考虑由下式给出的曲线 $c:[0,1]\to\mathbb{R}^2\times S^1$:
\[
  c(t):=
  \begin{cases}
   (x_0,y_0,3(\theta_L-\theta_0)t+\theta_0), & t\in[0,\tfrac13],\\
   (x_0+(3t-1)(x_1-x_0),\,y_0+(3t-1)(y_1-y_0),\,\theta_L), & t\in[\tfrac13,\tfrac23],\\
   (x_1,y_1,(3t-2)(\theta_1-\theta_L)+\theta_L), & t\in[\tfrac23,1],
  \end{cases}
\]
其中, 若 $(x_0,y_0)\ne(x_1,y_1)$, 则
\[
  (\cos\theta_L,\sin\theta_L)=\frac{(x_1-x_0,y_1-y_0)}{\sqrt{(x_1-x_0)^2+(y_1-y_0)^2}},
\]
否则 $\theta_L=\theta_0$. 显然 $c$ 连续且分段光滑, $c(0)=p$, $c(1)=q$. 此外,
\[
  \dot c(t)=
  \begin{cases}
    3(\theta_L-\theta_0)Y, & t\in(0,\tfrac13),\\
    3\sqrt{(x_1-x_0)^2+(y_1-y_0)^2}\,X, & t\in(\tfrac13,\tfrac23),\\
    3(\theta_1-\theta_L)Y, & t\in(\tfrac23,1)
  \end{cases}
\]
且
\[
  X=\cos\theta_L\frac{\partial}{\partial x}+\sin\theta_L\frac{\partial}{\partial y}
  \quad\text{和}\quad Y=\frac{\partial}{\partial\theta},
\]
所以 $c$ 与 $\Sigma$ 相容. 因此, 从 $p$ 出发经相容曲线可达的点的集合是 $\mathbb{R}^2\times S^1$, 所以 $\Sigma$ 不可能可积 (若 $\Sigma$ 可积, 则任一相容曲线都将局限于某一积分子流形).
\item 因为动能由
\[
  K=\frac M2\bigl((v^x)^2+(v^y)^2\bigr)+\frac I2(v^\theta)^2
\]
给出, 有
\begin{align*}
  \frac{\partial K}{\partial x}&=\frac{\partial K}{\partial y}=\frac{\partial K}{\partial\theta}=0,\\
  \frac{\partial K}{\partial v^x}&=Mv^x,\quad
  \frac{\partial K}{\partial v^y}=Mv^y,\quad
  \frac{\partial K}{\partial v^\theta}=Iv^\theta,
\end{align*}
因而
\[
  \mu\left(\frac{D\dot c}{\mathrm{d}t}\right)
  =M\ddot{x}\,\mathrm{d}x+M\ddot y\,\mathrm{d}y+I\ddot\theta\,\mathrm{d}\theta.
\]
此外, 因为约束反力是理想的, 有
\[
  \mathcal R=\lambda\omega=-\lambda\sin\theta\,\mathrm{d}x+\lambda\cos\theta\,\mathrm{d}y.
\]
于是冰刀的运动由运动方程
\[
  \mu\left(\frac{D\dot c}{\mathrm{d}t}\right)=\mathcal R(\dot c)
\]
的解给出, 它还须满足约束: $(\dot x,\dot y)$ 与 $(\cos\theta,\sin\theta)$ 成比例. 即它是常微分方程组
\[
  \begin{cases}
   M\ddot{x}=-\lambda\sin\theta,\\
   M\ddot y=\lambda\cos\theta,\\
   \ddot\theta=0,\\
   \dot y\cos\theta-\dot x\sin\theta=0
  \end{cases}
\]
的解. 因此, 对某个常数 $k\in\mathbb{R}$, 有 $\theta(t)=\theta_0+kt$.

若 $k\ne0$, 对 $\dot y\cos\theta-\dot x\sin\theta=0$ 求导得
\begin{align*}
  \ddot y=\ddot{x}\tan\theta+\frac{k}{\cos^2\theta}\dot x
  &\iff\frac{\lambda}{M}\cos\theta=\ddot{x}\tan\theta+\frac{k}{\cos^2\theta}\dot x\\
  &\iff-\frac{\cos\theta}{\sin\theta}\ddot{x}=\ddot{x}\tan\theta+\frac{k}{\cos^2\theta}\dot x\\
  &\iff\frac{\ddot{x}}{\dot x}=-k\tan\theta.
\end{align*}
于是 $\log|\dot x|=\log|\cos\theta|+\text{常数}$, 得到
\[
  \begin{cases}
   x(t)=r\sin(\theta_0+kt)+A_0,\\
   y(t)=-r\cos(\theta_0+kt)+B_0,\\
   \theta(t)=\theta_0+kt,\\
   \lambda=Mk^2r,
  \end{cases}
\]
其中 $r,A_0,B_0,\theta_0,k$ 是积分常数. 注意
\[
  (x(t)-A_0)^2+(y(t)-B_0)^2=r^2,
\]
所以 $c(t)$ 以恒定速率 $|kr|$ 描出以 $(A_0,B_0)\in\mathbb{R}^2$ 为圆心, 以 $|r|$ 为半径的圆. 可以将约束反力解释为不允许冰刀横向滑动的摩擦力, 它迫使轨迹弯曲.

若 $k=0$, 对 $\dot y\cos\theta-\dot x\sin\theta=0$ 求导得
\[
  \ddot y\cos\theta_0-\ddot{x}\sin\theta_0=0\iff\frac{\lambda}{M}=0,
\]
因此
\[
  \begin{cases}
   x(t)=l\cos\theta_0\,t+x_0,\\
   y(t)=l\sin\theta_0\,t+y_0,\\
   \theta(t)=\theta_0,\\
   \lambda=0,
  \end{cases}
\]
其中 $x_0,y_0,\theta_0,l$ 是积分常数. 注意, 这时 $c(t)$ 以恒定速率 $|l|$ 描出经过 $(x_0,y_0)$ 且斜率为 $\tan\theta_0$ 的直线. 因为冰刀沿自身长度方向滑动, 这时约束反力为零.
\item 现在运动方程为
\begin{align*}
  &\mu\left(\frac{D\dot c}{\mathrm{d}t}\right)=-\mathrm{d}U+\mathcal R(\dot c)\\
  &\quad\iff M\ddot{x}\,\mathrm{d}x+M\ddot y\,\mathrm{d}y+I\ddot\theta\,\mathrm{d}\theta
       =-Mg\sin\alpha\,\mathrm{d}x-\lambda\sin\theta\,\mathrm{d}x+\lambda\cos\theta\,\mathrm{d}y.
\end{align*}
于是冰刀的运动由该方程的一个同时满足约束方程的解给出, 即下列常微分方程组的一个解:
\[
  \begin{cases}
   M\ddot{x}=-Mg\sin\alpha-\lambda\sin\theta,\\
   M\ddot y=\lambda\cos\theta,\\
   \ddot\theta=0,\\
   \dot y\cos\theta-\dot x\sin\theta=0.
  \end{cases}
\]
因此, 对某个常数 $k\in\mathbb{R}$, 有 $\theta(t)=\theta_0+kt$. 若 $k\ne0$, 对 $\dot y\cos\theta-\dot x\sin\theta=0$ 求导得
\[
  \ddot y=\ddot{x}\tan\theta+\frac{k}{\cos^2\theta}\dot x
  \iff\frac{\lambda}{M}\cos\theta=\ddot{x}\tan\theta+\frac{k}{\cos^2\theta}\dot x.
\]
由常微分方程组的第一个方程得到
\[
  \frac{\lambda}{M}=-g\frac{\sin\alpha}{\sin\theta}-\frac1{\sin\theta}\ddot{x},
\]
所以代入上式得
\[
  -\frac{\cos\theta}{\sin\theta}(g\sin\alpha+\ddot{x})
     =\ddot{x}\tan\theta+\frac{k}{\cos^2\theta}\dot x,
\]
这蕴含
\[
  \frac{\ddot{x}}{\cos\theta}+k\frac{\sin\theta}{\cos^2\theta}\dot x
  =-g\sin\alpha\cos\theta
\]
或者等价地,
\[
  \frac{\mathrm{d}}{\mathrm{d}t}\left(\frac{\dot x}{\cos\theta}\right)
  =-g\sin\alpha\cos\theta.
\]
因此对某个积分常数 $l\in\mathbb{R}$, 有
\[
  \dot x=-\frac gk\sin\alpha\sin\theta\cos\theta+l\cos\theta
\]
和
\[
  \dot y=-\frac gk\sin\alpha\sin^2\theta+l\sin\theta.
\]
对最后这个方程求导得
\[
  \ddot y=(kl-2g\sin\alpha\sin\theta)\cos\theta,
\]
所以
\[
  \lambda=M(kl-2g\sin\alpha\sin\theta).
\]
于是得到
\[
  \begin{cases}
   x(t)=\dfrac{g}{4k^2}\sin\alpha\cos(2(\theta_0+kt))+\dfrac lk\sin(\theta_0+kt)+A_0,\\[4pt]
   y(t)=-\dfrac{g}{2k}\sin\alpha\left(t-\dfrac1{2k}\sin(2(\theta+kt))\right)
        -\dfrac lk\cos(\theta_0+kt)+B_0,\\[4pt]
   \theta(t)=\theta_0+kt,\\
   \lambda=M(kl-2g\sin\alpha\sin(\theta_0+kt)),
  \end{cases}
\]
其中 $\theta_0,k,l,A_0,B_0$ 是积分常数. 值得注意的是, 与可能的预期相反, $x(t)$ 保持有界, 而 $y(t)$ 线性增长.

若 $k=0$, 再次对 $\dot y\cos\theta-\dot x\sin\theta=0$ 求导, 得
\[
  \begin{cases}
   x(t)=-\dfrac g2\sin\alpha\cos^2\theta_0\,t^2+l\cos\theta_0\,t+x_0,\\[4pt]
   y(t)=-\dfrac g2\sin\alpha\sin\theta_0\cos\theta_0\,t^2+l\sin\theta_0\,t+y_0,\\[4pt]
   \theta(t)=\theta_0,\\
   \lambda=-Mg\sin\alpha\sin\theta_0,
  \end{cases}
\]
其中 $\theta_0,l,x_0,y_0$ 是积分常数. 如所预期, 这种情形下的运动是匀加速的, 加速度为 $g\sin\alpha\cos\theta_0$.
\end{enumerate}
\end{description}
